\chapter{Resultados}
\label{cap:resultados}

Esta se\c{c}\~ao apresenta os resultados preliminares do diagn\'ostico de
barreiras informacionais e d\'uvidas recorrentes do investidor de perfil
popular, executado na fase de Diagn\'ostico do cronograma da pesquisa
(Anexo~\ref{ane:cronograma}, ago.--out.\ 2026). Os achados constituem
resultados parciais, sujeitos a aprofundamento; as se\c{c}\~oes 5.2 a 5.5
permanecem pendentes, por dependerem de etapas emp\'iricas futuras
(engenharia dos m\'odulos Python, constru\c{c}\~ao da base RAG,
backtesting e avalia\c{c}\~ao t\'ecnica).

\section{Mapeamento das Barreiras Informacionais}
\label{sec:mapeamento-barreiras}

A estrat\'egia de coleta combinou tr\^es camadas de evid\^encia, da mais
robusta estatisticamente \`a mais ilustrativa, sintetizadas no
Quadro~\ref{qua:sintese-evidencias}.

\begin{quadro}[htbp]
  \centering
  \caption{S\'intese das Evid\^encias sobre Barreiras Informacionais}
  \label{qua:sintese-evidencias}
  \small
  \begin{tabular}{p{2.2cm}p{3.6cm}p{3.6cm}p{4.2cm}}
    \toprule
    \textbf{Camada} & \textbf{Fonte} & \textbf{Papel} & \textbf{Tipo de evid\^encia} \\
    \midrule
    1 -- Principal & Pesquisa de Letramento Financeiro (BCB/FGC, 2023) & Quantificar a barreira de forma estatisticamente representativa & Survey nacional, N=2.000, IC 95\% \\
    2 -- Complementar & Raio X do Investidor Brasileiro (ANBIMA, 2025/2026) & Perfis comportamentais e tend\^encias populacionais & Survey anual (Datafolha), N$\approx$5.800 \\
    3 -- Ilustrativa & Reclame Aqui (7 bancos p\'ublicos, registros filtrados) & Ilustrar a manifesta\c{c}\~ao pr\'atica da barreira & An\'alise de conte\'udo qualitativa \autocite{bardin2011} \\
    \bottomrule
  \end{tabular}
  \fonte{Elaborado pelo autor (2026).}
\end{quadro}

\subsection{Camada 1 -- Pesquisa de Letramento Financeiro (BCB/FGC, 2023)}
\label{subsec:camada1}

A evid\^encia quantitativa central do diagn\'ostico prov\'em da Pesquisa
de Letramento Financeiro, conduzida pelo Banco Central do Brasil (BCB) em
parceria com o Fundo Garantidor de Cr\'editos (FGC), com metodologia
padronizada internacionalmente pelo Toolkit OCDE/INFE \autocite{bcbfgc2023}
(2.000 entrevistas, amostra probabil\'istica estratificada, intervalo de
confian\c{c}a de 95\%, margem de erro de 2,19\%). Os resultados indicam
que apenas 48,2\% da popula\c{c}\~ao brasileira j\'a ouviu falar de
``T\'itulos financeiros (como Tesouro Direto ou CDBs)'' -- a categoria de
produto financeiro com menor \'indice de conhecimento entre as testadas,
\`a frente apenas de conta de investimento (31,0\%) e produtos
classificados como sustent\'aveis (16,6\%). Entre os entrevistados que
pouparam nos \'ultimos doze meses, apenas 7,1\% aplicaram em Tesouro
Direto ou CDB, contra 43,0\% que utilizaram caderneta de poupan\c{c}a ou
conta corrente.

Mais relevante ainda para a arquitetura proposta \'e o resultado da
dimens\~ao de conhecimento financeiro mensurado -- isto \'e, aferido por
meio de quest\~oes de c\'alculo, e n\~ao por autodeclara\c{c}\~ao. Embora
83,9\% dos entrevistados demonstrem compreender a no\c{c}\~ao geral de
juros, apenas 14,3\% conseguiram calcular corretamente um problema de
juros simples, e apenas 46,6\% responderam corretamente sobre a
implica\c{c}\~ao de juros compostos. Esse descolamento entre compreender
o conceito e conseguir oper\'a-lo numericamente corrobora, com dados
prim\'arios e representatividade nacional, a justificativa cient\'ifica
apresentada na Introdu\c{c}\~ao para a segrega\c{c}\~ao entre
processamento num\'erico determin\'istico e gera\c{c}\~ao de linguagem
natural.

\subsection{Camada 2 -- Raio X do Investidor Brasileiro (ANBIMA)}
\label{subsec:camada2}

A 8\textsuperscript{a} edi\c{c}\~ao do Raio X do Investidor Brasileiro
\autocite{anbima2025} fundamenta a segmenta\c{c}\~ao comportamental em
quatro perfis j\'a utilizada para calibrar as personas sint\'eticas do
Anexo~\ref{ane:personas}. A 9\textsuperscript{a} edi\c{c}\~ao
\autocite{anbima2026} acrescenta um dado relevante ao diagn\'ostico: do
total de brasileiros que conseguiram poupar em 2025, apenas 12\%
efetivamente converteram a poupan\c{c}a em investimento, evidenciando o
funil poupan\c{c}a--investimento como gargalo estrutural que antecede a
pr\'opria decis\~ao sobre qual produto de renda fixa escolher.

\subsection{Camada 3 -- Reclame Aqui (Bancos P\'ublicos)}
\label{subsec:camada3}

Como camada ilustrativa, foi conduzida an\'alise de conte\'udo qualitativa
\autocite{bardin2011} sobre registros p\'ublicos do Reclame Aqui filtrados
pelo conjunto de sete bancos p\'ublicos comerciais de varejo do Brasil --
Banco do Brasil, Caixa Econ\^omica Federal, BRB, Banrisul, Banestes,
Banese e Banpar\'a -- e por vincula\c{c}\~ao a decis\~ao ativa de
investimento em CDB ou Tesouro Direto, excluindo-se registros relativos a
produtos de aplica\c{c}\~ao autom\'atica sobre saldo em conta corrente
(ex.: CDB Sal\'ario no BRB, CDB Autom\'atico no Banrisul), que n\~ao
configuram decis\~ao de investimento por parte do cliente. O uso do
Reclame Aqui como fonte de pesquisa acad\^emica encontra precedente
metodol\'ogico em \textcite{rimolimelo2018}. A amostra resultante,
reduzida e n\~ao probabil\'istica, n\~ao deve ser lida como
representativa -- seu papel \'e ilustrar, com casos concretos, como as
barreiras identificadas nas Camadas 1 e 2 se manifestam na jornada real
do investidor de perfil popular em bancos p\'ublicos. O conte\'udo
relevante concentrou-se em tr\^es das sete institui\c{c}\~oes (Banco do
Brasil, Caixa e BRB); Banestes, Banese e Banpar\'a n\~ao apresentaram
registros indexados publicamente sobre CDB ou Tesouro Direto no per\'iodo
consultado, e o Banrisul apresentou majoritariamente registros sobre seu
produto de aplica\c{c}\~ao autom\'atica, exclu\'idos pelo mesmo crit\'erio
aplicado ao CDB Sal\'ario. Emergiram cinco categorias recorrentes: (i)
confus\~ao sobre o processo de cadastro e liquida\c{c}\~ao via CBLC/B3
para operar Tesouro Direto pelo banco -- a categoria mais saturada,
presente de forma independente no Banco do Brasil e na Caixa, com relatos
recorrentes de mensagens como ``aguardando retorno da CBLC'' interpretadas
como falha do sistema; (ii) fric\c{c}\~ao e custos ocultos ao diversificar
investimentos por meio de corretora parceira (BRB); (iii) desconhecimento
sobre a renova\c{c}\~ao autom\'atica (\textit{rollover}) do CDB no
vencimento (BRB); (iv) falta de transpar\^encia sobre o status do
investimento no aplicativo (BRB); e (v) desconhecimento sobre janelas
operacionais de resgate (BRB, Caixa).

N\~ao foi localizado, na literatura revisada, artigo com revis\~ao por
pares que combine an\'alise de conte\'udo de reclama\c{c}\~oes, renda
fixa e barreiras do investidor de perfil popular, o que indica uma
lacuna que esta pesquisa contribui para preencher.

\textbf{S\'intese integradora.} As tr\^es camadas de evid\^encia, tomadas
em conjunto, permitem hierarquizar as barreiras informacionais
identificadas da mais estrutural \`a mais operacional. No topo est\'a
uma barreira de conhecimento b\'asico: 51,8\% da popula\c{c}\~ao brasileira
sequer reconhece os termos ``Tesouro Direto'' ou ``CDB'' (Camada 1),
o que precede logicamente qualquer d\'uvida mais espec\'ifica sobre esses
produtos. Em seguida, uma barreira de opera\c{c}\~ao num\'erica: mesmo
entre quem compreende a no\c{c}\~ao geral de juros, apenas 14,3\% sabe
calcul\'a-los corretamente (Camada 1) -- gargalo que a arquitetura
proposta endere\c{c}a diretamente ao segregar o c\'alculo determin\'istico
da gera\c{c}\~ao de linguagem. Uma terceira barreira \'e de acesso: do
p\'ublico que consegue poupar, apenas 12\% converte a poupan\c{c}a em
investimento (Camada 2), evidenciando atrito na transi\c{c}\~ao entre
poupar e investir que antecede a escolha do produto. Por fim, duas
barreiras operacionais espec\'ificas e recorrentes aparecem na Camada 3,
ambas observadas em mais de uma institui\c{c}\~ao: a confus\~ao sobre o
processo de cadastro e liquida\c{c}\~ao via CBLC/B3 para operar Tesouro
Direto pelo banco (Banco do Brasil e Caixa) e o desconhecimento sobre
janelas operacionais de resgate (BRB e Caixa) -- o que lhes confere
maior robustez do que as demais tr\^es categorias da mesma camada, cada
uma restrita a um \'unico registro no BRB. Essa hierarquia --
conhecimento b\'asico $\rightarrow$ opera\c{c}\~ao num\'erica
$\rightarrow$ acesso $\rightarrow$ fric\c{c}\~ao operacional
espec\'ifica -- fornece o insumo priorit\'ario para a
calibra\c{c}\~ao das personas sint\'eticas (Anexo~\ref{ane:personas}) e
para a defini\c{c}\~ao do escopo inicial da base de conhecimento
regulat\'oria do agente (objetivo espec\'ifico b).

\bigskip
\textit{As demais se\c{c}\~oes deste cap\'itulo permanecem pendentes, a
serem redigidas ap\'os a execu\c{c}\~ao emp\'irica prevista no cronograma
(Anexo~\ref{ane:cronograma}):}

% TODO: preencher apos a execucao empirica.
\begin{itemize}
  \item[] \textit{5.2. Arquitetura do agente e m\'odulos Python -- valida\c{c}\~ao t\'ecnica (objetivos b e c) [PENDENTE -- base legal das tabelas tribut\'arias do RFL j\'a levantada em bibliografia/base-legal-rfl.md]}
  \item[] \textit{5.3. Base de conhecimento regulat\'oria RAG -- recall@k e prioriza\c{c}\~ao de norma vigente (objetivo d) [PENDENTE]}
  \item[] \textit{5.4. Resultados do backtesting de 12 meses (jun/25--mai/26) [PENDENTE]}
  \item[] \textit{5.5. Avalia\c{c}\~ao t\'ecnica automatizada (objetivo e) [PENDENTE]}
\end{itemize}
